%% kapa: exam=2 round=37 year=2026 period=1 subject=appraisal_practice subject_ko=감정평가실무 seq=2 number=2 points=30
%% source: Q-Net 감정평가사 기출문제 (cst00309) artlSeq=5268128 file=raw/1008/5268128/26년 제37회 감정평가사 2차 (1교시_감정평가실무).pdf
%% conversion: kordoc 4.24.51 -> tools/gen2.py (auto); review=proofread-sonnet; flags=none
\begin{longquestion}{【문제 2】}
2025. 1. 1.부터 주유소를 임차하여 운영중인 甲은 최근 관리청으로부터 토양오염 정기검사 결과 주유소 부지 지하에서 오염토양이 발견되어 토양정화를 통지 받았다. 이에 甲은 당해 주유소를 오랜기간 운영하여 토양오염 원인자로 추정되는 乙에게 토양정화에 소요되는 비용 및 토양정화 기간 동안 휴업으로 인한 피해액을 청구하기로 하고 청구금액을 확정하기 위해 감정평가사 丙에게 감정평가를 의뢰하였다. 제시된 자료와 관련법규 및 감정평가이론 등을 활용하여 다음 물음에 답하시오. (30점) \par
\medskip\noindent 물음 1) 토양정화 비용을 산정하시오. (15점) \par
\medskip\noindent 물음 2) 휴업으로 인한 피해액을 영업이익 손실액과 고정비 손실액으로 구분하여 산정하시오. (10점) \par
\medskip\noindent 물음 3) 의뢰인 甲이 감정평가사 丙에게 감정평가액을 5천만원 이상으로 요구할 경우, 감정평가사의 성실의무 등과 관련하여 약술하시오. (5점) \par
< 자료 1 > 토양오염 현황 \par
토양관련 전문기관에서 통보한 토양오염정밀조사보고서에 따르면, 오염물질은 주로 휘발유 성분이었고, 오염원인은 과거 유류 취급시설로부터의 유출로 인한 수평ㆍ수직 확산이 일어난 것으로 판단하였음 \par
< 자료 2 > 토양오염 수평 현황 \par
{\small\centering
\begin{tabular}{|c|c|c|}
\hline
\makecell{가로(m)} & \makecell{세로(m)} & \makecell{면적(m\textsuperscript{2})} \\
\hline
\makecell{3} & \makecell{3} & \makecell{9} \\
\hline
\end{tabular}\par}

< 자료 3 > 토양오염 수직 현황 \par
{\small\centering
\begin{tabular}{|c|c|c|c|}
\hline
\makecell{깊이(m)} & \makecell{토양 등 종류} & \makecell{오염면적(m\textsuperscript{2})} & \makecell{오염부피(m\textsuperscript{3})} \\
\hline
\makecell{지표면 ∼ 0.5} & \makecell{콘크리트} & \makecell{-} & \makecell{-} \\
\hline
\makecell{0.5 ∼ 1.0} & \makecell{점토} & \makecell{-} & \makecell{-} \\
\hline
\makecell{1.0 ∼ 2.0} & \makecell{점토} & \makecell{-} & \makecell{-} \\
\hline
\makecell{2.0 ∼ 3.0} & \makecell{점토} & \makecell{9} & \makecell{9} \\
\hline
\makecell{3.0 ∼ 4.0} & \makecell{점토} & \makecell{9} & \makecell{9} \\
\hline
\makecell{4.0 ∼ 5.0} & \makecell{점토} & \makecell{-} & \makecell{-} \\
\hline
\makecell{5.0 ∼ 6.0} & \makecell{점토} & \makecell{-} & \makecell{-} \\
\hline
\multicolumn{2}{|c|}{\makecell{합계}} & \makecell{-} & \makecell{18} \\
\hline
\end{tabular}\par}

< 자료 4 > 지표면 콘크리트 자르기 및 깨기 \par
가로 3m, 세로 3m 토양오염 면적만 정밀 작업하고, 콘크리트 자르기 및 깨기는 콘크리트 두께 0.5m를 고려한 단가이며, 100,000원/m\textsuperscript{2}임 \par
< 자료 5 > 토질별 팽창률 및 수축률(다짐률) \par
{\small\centering
\begin{tabular}{|c|c|c|}
\hline
\makecell{토질 분류} & \makecell{팽창률} & \makecell{수축률(다짐률)} \\
\hline
\makecell{모래} & \makecell{1.10} & \makecell{0.90} \\
\hline
\makecell{토사} & \makecell{1.20} & \makecell{0.85} \\
\hline
\makecell{점토} & \makecell{1.30} & \makecell{0.80} \\
\hline
\makecell{암석} & \makecell{1.50} & \makecell{1.20} \\
\hline
\end{tabular}\par}

< 자료 6 > 파쇄된 콘크리트 상차 및 폐콘크리트 처리비용 \par
파쇄된 콘크리트의 부피는 암석의 팽창률과 유사하며 폐콘크리트 처리비용은 팽창된 상태 기준으로 100,000원/m\textsuperscript{3}임 \par
< 자료 7 > 점토 굴착비용 \par
지표면 콘크리트 제거 후 점토 굴착비용은 오염되지 않은 부분과 오염된 부분이 동일하며, 굴착비용은 다짐된 점토 기준으로 20,000원/m\textsuperscript{3}임 \par
< 자료 8 > 덤프트럭 상차비용 \par
오염토양을 토양정화업자에게 보내기 위한 상차비용은 팽창된 점토 기준 10,000원/m\textsuperscript{3}이며, 오염되지 않은 점토는 청토(이미 오염정화된 점토, 이하 ‘청토’라 함)와 함께 당일 복토하고 다지므로 별도 비용이 발생하지 않음 \par
< 자료 9 > 덤프트럭 종류별 적재량 및 1회당 운반비용 \par
{\small\centering
\begin{tabular}{|c|c|c|}
\hline
\makecell{덤프트럭 종류} & \makecell{최대 적재부피(점토)} & \makecell{1회당 운반비용} \\
\hline
\makecell{5톤 덤프트럭} & \makecell{4m\textsuperscript{3}} & \makecell{200,000원} \\
\hline
\makecell{15톤 덤프트럭} & \makecell{10m\textsuperscript{3}} & \makecell{250,000원} \\
\hline
\makecell{25톤 덤프트럭} & \makecell{15m\textsuperscript{3}} & \makecell{300,000원} \\
\hline
\end{tabular}\par}

< 자료 10 > 오염토양 정화비용 \par
오염토양 정화비용은 토양정화 검증비용을 포함하여 팽창된 점토 기준 200,000원/m\textsuperscript{3} 임. 다만, 휴업으로 인한 피해액을 줄이기 위해 본건 오염토양을 오염토양정화처리장으로 반출하고, 기 정화완료된 청토를 대신 반입하여 오염토양정화하는 방법을 사용하며, 적법한 방법인 것으로 추정함 \par
< 자료 11 > 오염토양 정화된 점토 반입 및 다짐 비용 \par
정화처리된 오염토양의 반입비용은 토양정화업자가 부담하며, 다짐은 청토와 오염되지 않은 점토를 순차로 다짐. 다짐이 필요한 토량은 수축률(다짐률, 0.80)을 고려하여 산정하고, 다짐비용은 팽창된 점토기준 50,000원/m\textsuperscript{3}임. 점토의 팽창률(1.30)과 수축률(다짐률, 0.80)의 차이에 따른 잔여 점토의 처리비용은 다짐비용에 포함되어 고려하지 않음 \par
< 자료 12 > 레미콘 가격 단가표 \par
아래 표는 6m\textsuperscript{3} 용량 레미콘트럭 1대일 경우 단가이며, 주문량이 5m\textsuperscript{3} 미만일 경우에는 1회당 운반비 50,000원이 추가됨 \par
{\small\centering
\begin{tabular}{|c|c|c|}
\hline
\makecell{규격} & \makecell{단가/m\textsuperscript{3}} & \makecell{주요 용도} \\
\hline
\makecell{25-18-120} & \makecell{90,000원} & \makecell{버림 콘크리트} \\
\hline
\makecell{25-21-120} & \makecell{100,000원} & \makecell{주유소 바닥} \\
\hline
\makecell{25-24-150} & \makecell{110,000원} & \makecell{일반 건축물 골조} \\
\hline
\makecell{25-27-150} & \makecell{120,000원} & \makecell{고강도 구조물} \\
\hline
\end{tabular}\par}

< 자료 13 > 레미콘 타설 비용 \par
본건의 레미콘 타설에 필요한 비용은 철근, 철근 조립 인건비, 레미콘 타설 인건비를 포함하여 1,500,000원임 \par
< 자료 14 > 토양정화 작업기간 \par
오염토양 정화작업 기간내 당해 주유소는 차량 출입이 불가능하여 휴업이 필요한 것으로 판단됨. 휴업기간은 1개월이 소요될 것으로 추정되며, 피해액 산정은 월 단위를 기준으로 함 \par
< 자료 15 > 손익계산서(2025. 1. 1. ∼ 2025. 12. 31.) \par
{\small\centering
\begin{tabular}{|c|c|c|c|}
\hline
\makecell{과목} & \makecell{금액(원)} & \makecell{금액(원)} & \makecell{비고} \\
\hline
\makecell{Ⅰ.매출액} &  & \makecell{1,200,000,000} &  \\
\hline
\makecell{상품매출} & \makecell{1,200,000,000} &  &  \\
\hline
\makecell{Ⅱ.매출원가} &  & \makecell{900,000,000} &  \\
\hline
\makecell{상품매출원가} & \makecell{900,000,000} &  &  \\
\hline
\makecell{Ⅲ.매출총이익} &  & \makecell{300,000,000} &  \\
\hline
\makecell{Ⅳ.판매비와 관리비} &  & \makecell{231,500,000} &  \\
\hline
\makecell{임직원급여} & \makecell{100,000,000} &  & \makecell{자가노력비 포함} \\
\hline
\makecell{일용급여} & \makecell{10,000,000} &  & \makecell{직원휴가시 대리근무} \\
\hline
\makecell{퇴직급여} & \makecell{8,000,000} &  & \makecell{직원 퇴직금} \\
\hline
\makecell{복리후생비} & \makecell{15,000,000} &  & \makecell{직원 근무중 식사비} \\
\hline
\makecell{여비교통비} & \makecell{500,000} &  & \makecell{직원 업무관련 교통비} \\
\hline
\makecell{임차료} & \makecell{60,000,000} &  & \makecell{주유소 임차료} \\
\hline
\makecell{전력비} & \makecell{12,000,000} &  & \makecell{전기요금} \\
\hline
\makecell{보험료} & \makecell{2,000,000} &  & \makecell{자산 관련 보험료} \\
\hline
\makecell{감가상각비} & \makecell{4,000,000} &  &  \\
\hline
\makecell{지급수수료} & \makecell{20,000,000} &  & \makecell{신용카드 수수료} \\
\hline
\makecell{Ⅴ.영업이익} &  & \makecell{68,500,000} &  \\
\hline
\makecell{Ⅵ.영업외수익} &  & \makecell{300,000} &  \\
\hline
\makecell{잡수익} & \makecell{300,000} &  &  \\
\hline
\makecell{Ⅶ.영업외비용} &  & \makecell{13,000,000} &  \\
\hline
\makecell{이자비용} & \makecell{13,000,000} &  &  \\
\hline
\makecell{Ⅷ.당기순이익} &  & \makecell{55,800,000} &  \\
\hline
\end{tabular}\par}

< 자료 16 > 기타 사항 \par
1. 기준시점: 2026. 7. 4. \par
2. 피해액 산정을 위한 손익계산서는 주어진 자료를 사용할 것 \par
3. 산정가액은 백원단위에서 반올림하여 천원단위까지 표시함 \par
4. 주어진 자료 이외의 사항은 고려하지 아니함 \par
\end{longquestion}
